%============================================================
\chapter{Bai Hoc: Trach Nhiem Voi Cuoc Song (Responsibility for Life)}
\begin{center}
  \textit{\color{truyenblue}Gia cua su vi dai la trach nhiem; hay so huu qua quyet dinh, so huu cuoc doi.}\\
  \textit{(The price of greatness is responsibility; own your choices, own your life.)}\\[4pt]
  \small\color{truyengold}--- Thien nhien day ta ---
\end{center}
\ngancach

\section{Soi Bao Ve Lanh Dia}
\begin{truyen}{Soi Bao Ve Lanh Dia}{Con Soi}
\chuhoa{B}{}B ay soi co lanh dia va chung bao ve no khoi cac bay soi khac hoat dong xam pham.\\
\textit{(A wolf pack has territory and they protect it from other wolf packs infringing.)}

Nhung khi chung san qua muc, lanh dia suy kiet, bay soi nhan ra va tu han che san.\\
\textit{(But when they over-hunt, territory depletes, the pack notices and self-limits hunting voluntarily.)}

Khong co nha nuoc nao bat chung; chinh ban nang va tri nho tap the thuc day.\\
\textit{(No state forces them; their own instinct and collective memory drives this.)}

Khi deer population tang lai, bay soi dan dan mo rong sang muc do hop ly de tiep tuc.\\
\textit{(When deer populations recover, the pack gradually expands hunting to reasonable levels again.)}

Trach nhiem voi lanh dia va nguon thuc pham la su song con co the duy tri lau dai.\\
\textit{(Responsibility for territory and food sources is how survival can be maintained for the long term.)}

\end{truyen}
\begin{baihoc}
  Bao nhieu ban co the lay tu moi truong va bao nhieu ban se de lai cho tuong lai la cau hoi cua trach nhiem.\\
  \textit{(How much you take from an environment and how much you leave for the future is the question of responsibility.)}
\end{baihoc}

\section{Ong Chua Dieu Phoi Dan Ong}
\begin{truyen}{Ong Chua Dieu Phoi Dan Ong}{Ong Chua}
\chuhoa{O}{}O ng chua khong pha di sang de ngu hay nghi ngoi; ngoai viec san xuat trung nghia vu cua no ro rang.\\
\textit{(The queen bee does not wander outside to sleep or rest; besides laying eggs her duty is clear.)}

San xuat den 2000 trung moi ngay suot bai muoi ngay song, con thang, ngay nay qua ngay khac.\\
\textit{(Laying up to 2000 eggs per day throughout her sixty-day lifespan, month after month, day after day.)}

Neu ba ngung mot ngay, ca dan ong mat track phat trien dan so va co the suy sup trong vai thang.\\
\textit{(If she stops even one day, the whole colony loses population growth track and may collapse in months.)}

Ba khong bao gio can ai nhac nho bong bo; trach nhiem cua ba la ban nang cot loi.\\
\textit{(She never needs anyone reminding her; her responsibility is her core instinct.)}

Nguoi lanh dao khong can nhac nho moi lam viec; viec lam viec la chinh ho, khong can ung thuan.\\
\textit{(A leader doesn't need reminders to work; working IS who they are, requiring no encouragement.)}

\end{truyen}
\begin{baihoc}
  Trach nhiem cao nhat la khi ban lam viec vi tinh chat cong viec chinh no chu khong vi phan thuong hay gio giac.\\
  \textit{(The highest responsibility is when you work because of the very nature of the work, not for reward or schedule.)}
\end{baihoc}

\section{Ca He Bao Ve San Ho}
\begin{truyen}{Ca He Bao Ve San Ho}{Ca He}
\chuhoa{C}{}C a he nho bi nhan nhang soi bong den tren rao san ho noi no sinh song va con cai.\\
\textit{(The small clownfish is seen patrolling the sea anemone area where it lives and raises its young.)}

No cham soc kem thuc an cho so goi bien de giu suc khoe khoang san ho.\\
\textit{(It brings food to the sea anemone to maintain health of the coral neighborhood.)}

No se tan cong bat ky con ca nao lon hon no gan 20 lan ma dam den lau vung tam cua no.\\
\textit{(It will attack any fish up to 20 times its size that dares to come near its protected zone.)}

No chi co mot cai ngoi, mot cai khu vuc trach nhiem, va no bien no thanh ca mot the gioi de bao ve.\\
\textit{(It has only one home, one area of responsibility, and it turns it into an entire world worth protecting.)}

Trach nhiem khong nhat thiet phai lon moi co y nghia; bao ve tot nhung gi nho nhat cung la vi dai.\\
\textit{(Responsibility need not be large to be meaningful; protecting well the smallest things is also greatness.)}

\end{truyen}
\begin{baihoc}
  Hay nhan trach nhiem voi moi truong nho cua ban truoc; khi ban lam tot do, ban moi co the mo rong.\\
  \textit{(Take responsibility for your small environment first; when you do that well, you can expand.)}
\end{baihoc}

\section{Voi Nho De Lai Canh Dong}
\begin{truyen}{Voi Nho De Lai Canh Dong}{Bay Voi}
\chuhoa{K}{}K hi bay voi di qua canh rung, chung de lai nhung khoang trong boi voi o la cac cay xa.\\
\textit{(When elephant herds pass through a forest, they leave open spaces from pulling out trees for food.)}

Nhung khoang trong nay cho anh sang vao dat de hat giong nay mam va cac loai co thi de bay.\\
\textit{(These open spaces let sunlight reach the soil for seeds to sprout and short-grass species to seed.)}

Que chim kien va soc co dan theo vao o trong khoang trong do va bat dau tao nen he sinh thai rieng.\\
\textit{(Bird, ant and squirrel colonies follow into those clearings and begin forming their own ecosystem.)}

Sac de lai tren canh rung cua bay voi khong phai tan pha ma la tai tao da dang sinh hoc.\\
\textit{(Marks left on the forest by elephant herds are not destruction but biodiversity recreation.)}

Trach nhiem voi moi truong doi khi co nghia la tao ra nhung thay doi co chon ve vi long tot.\\
\textit{(Responsibility to the environment sometimes means creating intentional changes for the greater good.)}

\end{truyen}
\begin{baihoc}
  Trach nhiem voi the gioi tu nhien la hieu rang tac dong cua ta co the la qua tang hoac la chan thuong.\\
  \textit{(Responsibility to the natural world is understanding that our impact can be either a gift or an injury.)}
\end{baihoc}

\section{Nhay San Lac Khong Ro Nguoi}
\begin{truyen}{Nhay San Lac Khong Ro Nguoi}{Nhay San}
\chuhoa{M}{}M ot loai nhay san kinh nghiem biet rang mot khi da bat dau tan cong phuc kich thi phai hoan thanh.\\
\textit{(An experienced jumping spider knows that once a pounce attack begins it must be completed.)}

Dung lai giua chung chac chan se giup con moi thoat va tiep lat net mang nhot lua che ban than.\\
\textit{(Stopping midway would certainly let the prey escape and expose its web-concealment camouflage.)}

No la luon chiu het hau qua cua quyet dinh tan cong du co vang hay thua cuoc.\\
\textit{(It always bears the full consequences of the attack decision whether winning or losing.)}

Khong co chuyen nhay san dat len nua roi lui lai vi so mat mat chac tu moi truong khac.\\
\textit{(There is no such thing as a jumping spider beginning its leap then retreating out of fear of loss.)}

Trach nhiem la biet rang moi quyet dinh den hau qua va ta can san sang chiu hau qua do.\\
\textit{(Responsibility is knowing every decision has consequences and we must be ready to bear those consequences.)}

\end{truyen}
\begin{baihoc}
  Truoc khi do quyet dinh, hay suy nghi ky; nhung sau khi da quyet dinh, hay chiu trach nhiem den cung.\\
  \textit{(Before making a decision, think carefully; but once decided, bear responsibility to the very end.)}
\end{baihoc}

\section{Rat Binh Dan Giua Bieu Ro}
\begin{truyen}{Rat Binh Dan Giua Bieu Ro}{Con Rat}
\chuhoa{T}{}T rong nghien cuu noi tieng, con rat duoc dat trong hop co nut bam giai phong dong loai bi chum.\\
\textit{(In a famous study, a rat was placed in a box with a button that released a caged companion.)}

Con rat khong huong loi gi tu viec nhan ban nhu; no chi nhan ban boi no cam thay phai lam.\\
\textit{(The rat gained nothing from pressing the button; it only pressed because it felt it had to.)}

Chay qua den phan thuong thuc an de nhan ban truoc, roi moi an thuc an sau khi ban da tu do.\\
\textit{(It ran past food rewards to press the release first, then ate food after its companion was free.)}

Nhieu con rat con chia phan thuc an cua minh cho dong loai vua duoc giai thoat ra.\\
\textit{(Many rats even shared their own food portions with the companion just freed.)}

Trach nhiem voi dong loai, su dong cam, khong phai chi la cua con nguoi; no la dieu toan cau.\\
\textit{(Responsibility to fellow creatures, empathy, is not uniquely human; it is universal.)}

\end{truyen}
\begin{baihoc}
  Neu con rat biet trach nhiem voi dong loai thi con nguoi cang phai biet hon the rat nhieu.\\
  \textit{(If a rat knows responsibility to its kind then humans certainly must know it far more.)}
\end{baihoc}

\section{Ca Ham Co Den Lap Lua Cho Dan}
\begin{truyen}{Ca Ham Co Den Lap Lua Cho Dan}{Ca Ham}
\chuhoa{C}{}C a ham co den dung kha nang phat sang cua minh de dan dan ca nho theo sau.\\
\textit{(Anglerfish use their bioluminescence ability to lead small fish to follow behind.)}

Cac loai ca phiep tang theo lan anh sang de di theo cuoi cung bi an truoc khi hieu ra be.\\
\textit{(Naive fish are drawn to follow the light ultimately being eaten before realizing the trap.)}

Nhung mot so nha nghien cuu tim thay loai ca co den cung dung den trong muc dich cung sinh.\\
\textit{(But some researchers found anglerfish also use their light for mutualistic purposes.)}

Trong vung nuoc toi den cach bo bien 2000 met, hau nhu khong co loai nao tu lam den.\\
\textit{(In pitch-dark zones 2,000 meters from shore almost no species can produce their own light.)}

Trach nhiem cua nguoi co tai nang rieng biet la su dung tai nang do cho dieu lon lao hơn luoi ich ca nhan.\\
\textit{(The responsibility of one with a unique gift is to use that gift for something greater than personal gain.)}

\end{truyen}
\begin{baihoc}
  Bat ky tai nang gi ban co cung dua theo trach nhiem su dung no mot cach co ich cho nguoi khac.\\
  \textit{(Whatever talent you possess carries the responsibility of using it usefully for others.)}
\end{baihoc}

\section{Khi Di Cung Voi Con Gia Benh}
\begin{truyen}{Khi Di Cung Voi Con Gia Benh}{Con Khi}
\chuhoa{K}{}K hi nghien cuu quan sat mot con khi gia bi gai tre be nan di cham, bay khi van di cung toc do do.\\
\textit{(When researchers observed an elderly monkey with a bamboo splinter injury walking slowly, the whole troop walked at its pace.)}

Bay khi hoan toan co kha nang bo sung song truoc nhung khong con nao lam vay.\\
\textit{(The troop was fully capable of leaving the injured one behind but not one did so.)}

Nhung con khoe manh khi phuc ben nguoi yeu, quan sat, san sang can thiep khi can.\\
\textit{(The healthy ones watched beside the weaker one, observing, ready to intervene when needed.)}

Khi khi gia chet di, ca bay ngoi im lang ben xac no trong vai gio truoc khi di.\\
\textit{(When the old monkey died, the whole troop sat quietly beside the body for hours before leaving.)}

Trach nhiem voi nguoi yeu duoi khong bao gio la ganh nang; do la dieu co nghia la cuoc doi.\\
\textit{(Responsibility for the weaker is never a burden; it is what gives life meaning.)}

\end{truyen}
\begin{baihoc}
  Toc do cua mot nhom duoc do bang toc do cham nhat cua ho; dung bo roi nguoi chay cham nhat.\\
  \textit{(A group's speed is measured by the speed of its slowest member; never leave behind the slowest runner.)}
\end{baihoc}

\section{Ran Hoi Khong Can Nhan Nhung Khong Het}
\begin{truyen}{Ran Hoi Khong Can Nhan Nhung Khong Het}{Ran Hoi}
\chuhoa{K}{}K hi ran hoi cat vao con moi, no khong pha co doc toi da ngay ca khi co the lam vay.\\
\textit{(When a king cobra bites prey, it does not inject maximum venom even when it could.)}

Cac nghien cuu ghi nhan con ran luon dieu chinh lieu luong doc phu hop voi kich co moi.\\
\textit{(Studies record that the snake always calibrates venom dosage appropriate to the prey's size.)}

Dieu nay bao ton nguon co doc quy gia de su dung khi thuc su can biet trong tuong lai.\\
\textit{(This conserves the precious venom supply for truly necessary future use.)}

No khong lang phi suc manh; no su dung dung du suc manh can thiet va khong hon nua.\\
\textit{(It does not waste strength; it uses exactly the strength needed and no more.)}

Trach nhiem voi suc manh minh co la khong su dung qua muc can thiet bao gio.\\
\textit{(Responsibility with the power you possess means never using more than what is needed.)}

\end{truyen}
\begin{baihoc}
  Nguoi manh trach nhiem nhat su dung it suc manh nhat; ho biet rang quyen luc lon nen duoc su dung can than.\\
  \textit{(The most responsible powerful person uses the least power; they know great power should be used sparingly.)}
\end{baihoc}

\section{Dan Ca Bao Ve Bien}
\begin{truyen}{Dan Ca Bao Ve Bien}{Dan Ca Nho}
\chuhoa{N}{}N hieu dan toc bien di bien bao nhieu doi van giu nguyen suc song phong khoang cua nguon ca.\\
\textit{(Many coastal communities have fished the sea for generations yet kept the fish populations thriving.)}

Ho co luat khong thanh van: khong san mua ca de, khong san ca con, khong dan pha rao san ho.\\
\textit{(They have unwritten laws: no fishing in spawning season, no catching juvenile fish, no dragging over coral reefs.)}

Luat nay khong den tu khoa hoc hai duong hoc; no den tu tri tue dieu nhai qua hang tram nam.\\
\textit{(These laws did not come from marine science; they came from fine-tuned wisdom across hundreds of years.)}

Moi nguoi dan bien biet rang neu ho lay qua nhieu hom nay, con cai ho se doi kho ngay mai.\\
\textit{(Every fisherman knows that if they take too much today, their children will go hungry tomorrow.)}

Trach nhiem voi nguon tai nguyen la loai tri tue lau doi nhat va quan trong nhat tren hanh tinh.\\
\textit{(Responsibility for natural resources is the oldest and most important wisdom on the planet.)}

\end{truyen}
\begin{baihoc}
  Hay nghi den nhung nguoi chua sinh ra khi ban dua ra quyet dinh ve tai nguyen ngay hom nay.\\
  \textit{(Think of the unborn when making decisions about resources today.)}
\end{baihoc}
